\documentclass[11pt]{article}
\usepackage[T1]{fontenc}
\usepackage{lmodern}
\usepackage[margin=1in]{geometry}
\usepackage{amsmath,amssymb,amsthm,booktabs,array}
\usepackage[colorlinks=true,urlcolor=blue,linkcolor=blue,citecolor=blue]{hyperref}
\usepackage{microtype}
\newtheorem{theorem}{Theorem}
\newtheorem{lemma}[theorem]{Lemma}
\newtheorem{proposition}[theorem]{Proposition}
\newcommand{\Z}{\mathbb Z}
\newcommand{\Q}{\mathbb Q}
\newcommand{\OK}{\mathcal O_K}
\newcommand{\zK}{\zeta_K}
\newcommand{\tht}{\Theta}
\newcommand{\norm}[1]{\lvert #1\rvert}
\title{An explicit tenth-discriminant bound\\Version 2: the refined cutoff}
\author{Research note prepared with Codex}
\date{Version 0.2.0 --- 20 October 2026}
\begin{document}
\maketitle
\begin{abstract}
Taking explicitly identified results of Banks
(\texttt{arXiv:2610.15483v1}) as external inputs, we obtain
\(q<5\cdot10^7\) for imaginary quadratic fields of class number one.
This self-contained cumulative version retains the original
\(q<10^{10}\) argument, proves the refinement, and includes the finite
exclusion that leaves exactly the nine familiar fundamental discriminants.
The refinement changes both the theta estimate and the parameters in
Banks's determinant argument. Neither version uses Heegner's methods
or Banks's separate small-conductor zero computation.
\end{abstract}

\section{What changed since version 1}\label{sec-changes}

Version 1 remains available unchanged, including its source, PDF,
programs, and recorded outputs. This version contains the whole argument
needed to read and reproduce the refined result.
\begin{itemize}
\item The analytic cutoff improves from \(10^{10}\) to
\(5\cdot10^7\), a factor of \(200\).
\item The original theta estimate \(1-\beta<10/\sqrt q\) is retained
in Section~\ref{sec-baseline}. The refinement proves
\(1-\beta<2/\sqrt q\) for \(q\ge49\,000\,000\).
\item The original argument uses Banks's stated Theorem~1.1. The
refinement uses Lemma~4.2, Lemma~5.1, and Corollary~6.3, including their
preceding construction and hypotheses, with \(H=10^9\) and \(N=q^{18}\).
Thus the change is not a substitution of a smaller threshold into
Theorem~1.1.
\item The two original integer sieve programs and their outputs are
preserved byte-for-byte. Their search below \(10^{10}\) also covers
the smaller remaining interval. A new exact rational certificate
checks the constants used in the refinement.
\end{itemize}

The repository owner reports a cross-check of the original mathematical
argument with a different frontier model. The new refinement has not
yet received that separate-model cross-check. Its numerical inequalities
have passed the exact arithmetic certificate. These are distinct
verification records, detailed in Section~\ref{sec-audit}.

\section{Results and external dependencies}

Throughout, \(D<0\) is a fundamental discriminant,
\(q=\norm D\), \(K=\Q(\sqrt D)\), and \(\chi_D\) is the primitive
quadratic character of conductor \(q\). All logarithms are natural.
The quantity \(q\) is the absolute field discriminant, not necessarily
the squarefree integer under the square root.

\begin{theorem}[Refined bound]\label{main}
Assume Lemma~4.2, Lemma~5.1, and Corollary~6.3 of Banks~\cite{banks},
including their preceding construction and hypotheses.
If \(h(K)=1\), then
\[
 q<5\cdot10^7.
\]
Together with the finite verification in Section~\ref{finite}, this gives
\[
 q\in\{3,4,7,8,11,19,43,67,163\}.
\]
Every discriminant on this list does have class number one.
\end{theorem}

For comparison, the original external input was Banks's Theorem~1.1:
\begin{equation}\label{banks}
 q\ge10^{10},\quad L(\beta,\chi)=0,\quad 0<\beta<1
 \quad\Longrightarrow\quad
 (1-\beta)\log q\ge\frac3{1000}
\end{equation}
for every primitive real character \(\chi\) of conductor \(q\).

\begin{proposition}[Original bound]\label{baseline}
Assuming \eqref{banks}, if \(h(K)=1\), then \(q<10^{10}\).
\end{proposition}

The external assumptions are stated explicitly because this note does
not independently prove Banks's theorem or the intermediate lemmas.
The contribution is their explicit application to the class-number-one
problem, with reproducible finite exclusion. No novelty claim is made
for the classical theta method, and no optimality claim is made for
the numerical cutoff.

\section{Original argument: the cutoff \texorpdfstring{\(10^{10}\)}{10 to the tenth power}}\label{sec-baseline}

\begin{lemma}\label{zero}
If \(q\ge10^{10}\) and \(h(K)=1\), then \(L(s,\chi_D)\) has a real zero
\[
 1-\frac{10}{\sqrt q}<\beta<1.
\]
\end{lemma}

\begin{proof}
Put
\[
 R=\sqrt q,\qquad y=\frac R2,\qquad
 \Lambda=y^{-1/2}\OK\subset\mathbb C\simeq\mathbb R^2.
\]
The embedded integer lattice \(\OK\) has covolume \(y\), so
\(\Lambda\) has covolume one. Define
\[
 \tht(t)=\sum_{\alpha\in\OK}
       \exp\left(-\frac{\pi t\norm{\alpha}^2}{y}\right)
       \qquad(t>0).
\]
For any planar lattice of covolume one, its Euclidean dual is a
rotation by \(90\) degrees of the original lattice. For example, if
the columns of \(A\) are an oriented basis and \(\det A=1\), then
\(A^{-T}=JAJ^{-1}\), where \(J\) is rotation by \(90\) degrees
and \(J^{-1}\Z^2=\Z^2\).
Consequently, Poisson summation for the Gaussian gives
\begin{equation}\label{thetafunctional}
 \tht(t)=t^{-1}\tht(1/t).
\end{equation}

Since \(q>4\), the units of \(\OK\) are exactly \(\{\pm1\}\).
Class number one implies that every nonzero integral ideal has
exactly two generators in \(\OK\). Thus, initially for \(\Re s>1\),
\[
 \sum_{\alpha\in\OK\setminus\{0\}}\norm{\alpha}^{-2s}
     =2\zK(s).
\]
Taking the Mellin transform and splitting its integral at \(t=1\),
then using \eqref{thetafunctional} in the part below \(1\), gives
\begin{equation}\label{mellin}
 2\Gamma(s)\left(\frac y\pi\right)^s\zK(s)
 =
 -\frac{1}{s(1-s)}
 +\int_1^\infty(\tht(t)-1)
       \left(t^{s-1}+t^{-s}\right)\,dt .
\end{equation}
The integral on the right converges locally uniformly for all complex
\(s\), since \(\tht(t)-1\) decays exponentially as \(t\to\infty\).
Hence \eqref{mellin} continues the initial identity meromorphically.
For clarity, the rational term is
\[
 \int_0^1(t^{s-2}-t^{s-1})\,dt
 =\frac{1}{s-1}-\frac1s=-\frac{1}{s(1-s)}.
\]

Set
\[
 \varepsilon=\frac{10}{R},\qquad s_0=1-\varepsilon,\qquad
 A=\frac y\pi=\frac{R}{2\pi}.
\]
Here \(R\ge10^5\), \(0<s_0<1\), and \(A>1\).
The two lattice points \(\alpha=1,-1\) alone give
\[
 \tht(t)-1\ge2e^{-\pi t/y}\ge\frac2e
       \qquad(1\le t\le A).
\]
Dropping the nonnegative \(t^{-s_0}\) term and the part of the integral
above \(A\), the integral in \eqref{mellin} is therefore at least
\[
 \frac2e\int_1^A t^{s_0-1}\,dt
    =\frac{2}{e\,s_0}(A^{s_0}-1).
\]
It suffices to prove
\begin{equation}\label{target}
 \frac{2\varepsilon}{e}(A^{1-\varepsilon}-1)>1.
\end{equation}

All estimates needed for \eqref{target} can be made rational.
The function \((\log R)/R\) decreases for \(R>e\), and
\(\log(10^5)<12\). Hence
\[
 0<\varepsilon\log A
 \le\frac{10\log R}{R}<\frac{3}{2500}.
\]
Using \(e^{-u}\ge1-u\),
\[
 A^{1-\varepsilon}
   = A e^{-\varepsilon\log A}
   >\frac{R}{2\pi}\frac{2497}{2500}.
\]
The elementary bounds \(\pi<22/7\) and \(2<e<3\) now imply
\begin{align*}
 \frac{2\varepsilon}{e}(A^{1-\varepsilon}-1)
 &>\frac{10}{\pi e}\frac{2497}{2500}-\frac{20}{eR}\\
 &>\frac{35}{33}\frac{2497}{2500}-\frac{1}{10000}\\
 &=\frac{31777}{30000}>1.
\end{align*}
For an entirely rational check of the logarithmic bound, use
\(e>8/3\) and \(8^{12}>10^5\,3^{12}\).
This proves \eqref{target}, and so \(\zK(s_0)>0\).

As \(s\to1^-\), the rational term in \eqref{mellin} tends to
\(-\infty\), while its integral stays finite. Thus \(\zK(s)<0\)
sufficiently close to \(1\) from the left.
Continuity gives a zero \(\beta\in(s_0,1)\).
Finally \(\zK(s)=\zeta(s)L(s,\chi_D)\), and \(\zeta(s)\ne0\) for
real \(0<s<1\). One elementary proof of the latter uses the positive
alternating series
\[
 \eta(s)=\sum_{n\ge1}(-1)^{n-1}n^{-s}>0,\qquad
 \zeta(s)=\frac{\eta(s)}{1-2^{1-s}}<0.
\]
Therefore \(L(\beta,\chi_D)=0\), as claimed.
\end{proof}

\begin{proof}[Proof of Proposition~\ref{baseline}]
If \(q\ge10^{10}\) and \(h(K)=1\), Lemma~\ref{zero} gives
\[
 (1-\beta)\log q < \frac{10\log q}{\sqrt q}.
\]
The function \((\log q)/\sqrt q\) decreases for \(q>e^2\).
Using \(\log(10^{10})<24\), we obtain
\[
 (1-\beta)\log q
 < \frac{10\cdot24}{10^5}
 =\frac{3}{1250}<\frac{3}{1000}.
\]
This contradicts \eqref{banks}.
\end{proof}

\section{Refinement A: a stronger zero-free estimate}\label{sec-refined-banks}

We first show that those lemmas imply the following numerical statement:
for every primitive real character of conductor \(q\ge 5\cdot10^7\),
every real zero \(0<\beta<1\) satisfies
\begin{equation}\label{lower}
             \delta:=(1-\beta)\log q>\frac{509}{100000}=0.00509.
\end{equation}
Put
\[
 H=10^9,\qquad N=q^{18},\qquad M=N^{4/3}=q^{24},
 \qquad L=\log q.
\]
Take \(\lambda=3/2\) in Banks's prime-bias input. Its constants satisfy
\[
 C_1=\frac{2e^{3/2}}9<1,\qquad
 C_2=\frac{e^{3/2}}4<\frac98.
\]
The hypotheses \(q\ge9\), \(319\le H\le N\), \(M\ge8\), and
\(H^{2/3}N^{4/3}\ge 2\cdot97^2\) all hold.
In particular, Banks's construction requires \(H\le N\), not \(H\le q\).
The latter was a sufficient condition used with his original parameters,
and is not needed here.

Use the nonzero determinant and the degree sums \(S_1,S_{23}\) in
Banks's notation. Lemma~4.2 gives
\[
 S_1\ge \frac{10^6}{4\cdot97^2}N^4M>25N^4M,\qquad
 \frac{S_{23}}{S_1}\le\frac{768\cdot97^2}{10^9}<\frac1{135}.
\]
Combining Lemma~5.1 with Corollary~6.3, dividing by \(S_1>0\),
and using \(\log N=\tfrac34\log M\), gives
\begin{align}
 \frac{11}{45}\log M
 &\le
 \delta\frac{(\log M)^2}{L}
 +\left(\frac98+\frac1{2H}+\frac{68}{135}\right)L
 +\log H+\frac2{25}+\frac{136}{135}\log8
 +\frac{2\log N}{25M}.
 \label{comparison}
\end{align}
Here \(11/45=1-(136/135)(3/4)\). Replacing \(C_1\) by \(1\)
increases the right side because \(\delta>0\).

The coefficient of \(L\) after the first term is less than \(163/100\).
Also
\[
 \log H<9(2.303),\qquad \log8<2.08,\qquad
 \frac{2\log N}{25M}<\frac1{1000},
\]
so the remaining terms on the right of \eqref{comparison} sum to
less than \(23\). For the last inequality, \(N\ge1000\) and the function
\((\log N)/N^{4/3}\) is decreasing there; use \(\log1000<7\).
Since \(\log M=24L\), \eqref{comparison} implies
\[
 576\delta>\frac{88}{15}-\frac{163}{100}-\frac{23}{L}.
\]
Finally \(L\ge\log(5\cdot10^7)>17.7\), and hence
\begin{equation}\label{rationaldelta}
 \delta>
 \frac{1}{576}\left(\frac{88}{15}-\frac{163}{100}
                         -\frac{230}{177}\right)
 >\frac{509}{100000}.
\end{equation}
This proves \eqref{lower}, without a computation of any Dirichlet zeros.

\section{Refinement B: a closer real zero}\label{sec-refined-theta}

\begin{lemma}\label{theta}
If \(q\ge49\,000\,000\) and the class number is one, then there is
a real zero \(0<\beta<1\) of \(L(s,\chi_{-q})\) such that
\[
                         1-\beta<\frac2{\sqrt q}.
\]
\end{lemma}

\begin{proof}
Use the theta function and Mellin identity \eqref{mellin} derived
in Section~\ref{sec-baseline}, with \(R=\sqrt q\) and \(y=R/2\).
That identity requires only \(q>4\) and class number one; the stronger
range restriction there was used for the subsequent numerical estimate.
Here \(R\ge7000\).

Set
\[
 \varepsilon=\frac2R,\qquad s=1-\varepsilon,\qquad
 A=\frac{R}{2\pi},\qquad T=5A.
\]
Retain just the lattice points \(\alpha=\pm1,\ldots,\pm30\)
and just the \(t^{s-1}\) term over \(1\le t\le T\).
If \(I\) denotes the integral in \eqref{mellin}, then
\begin{align*}
 I
 &\ge 2T^{-\varepsilon}
        \sum_{n=1}^{30}\int_1^T e^{-n^2t/A}\,dt\\
 &=2A T^{-\varepsilon}
        \sum_{n=1}^{30}\frac{e^{-n^2/A}-e^{-5n^2}}{n^2}\\
 &\ge2A T^{-\varepsilon}
       \left(\sum_{n=1}^{30}\frac1{n^2}
              -\frac{30}{A}
              -\sum_{n=1}^{30}\frac{e^{-5n^2}}{n^2}\right).
\end{align*}
The last step uses \(e^{-u}\ge1-u\).
For explicit bounds, use
\[
 \sum_{n=1}^{30}\frac1{n^2}>\frac{1612}{1000},\qquad
 \sum_{n=1}^{30}\frac{e^{-5n^2}}{n^2}
 \le\frac{e^{-5}}{1-e^{-15}}<\frac7{1000}.
\]
Indeed \(n^2\ge1+3(n-1)\); and the second strict inequality follows
already from \(e>19/7\).
Since \(R\ge7000\), \(3<\pi<22/7\), and \((\log R)/R\) decreases,
\[
 \frac{30}{A}<\frac{1320}{49000},\qquad
 s\ge1-\frac2{7000},\qquad
 T^{-\varepsilon}\ge1-\varepsilon\log T
      >1-\frac{18}{7000}.
\]
For the last inequality, \(T<R\) and \(\log7000<9\).
It follows that
\begin{align*}
 \varepsilon sI
 &>
 \frac7{11}\left(1-\frac2{7000}\right)
            \left(1-\frac{18}{7000}\right)
 \left(\frac{1612}{1000}-\frac{1320}{49000}
                         -\frac7{1000}\right)\\
 &=\frac{37781022837}{37730000000}>1.
\end{align*}
Thus \(\zK(1-2/R)>0\) by \eqref{mellin}. The same identity shows
that \(\zK(s)\to-\infty\) as \(s\to1^-\), so continuity gives a
zero \(\beta\in(1-2/R,1)\).
Since \(\zK(s)=\zeta(s)L(s,\chi_{-q})\) and
\(\zeta(s)<0\) for real \(0<s<1\), this is a zero of the Dirichlet
\(L\)-function.
\end{proof}

\section{Conclusion of the refined bound}

If \(q\ge5\cdot10^7\) had class number one, Lemma~\ref{theta} would give
\[
 (1-\beta)\log q<\frac{2\log q}{\sqrt q}.
\]
The function on the right decreases for \(q>e^2\).
Using \(\log(5\cdot10^7)<17.73\) and
\(\sqrt{5\cdot10^7}>7071\), we therefore get
\[
 (1-\beta)\log q
 <\frac{2(17.73)}{7071}
 <\frac{502}{100000}=0.00502.
\]
This contradicts \eqref{lower} and proves Theorem~\ref{main}.

\section{Finite exclusion: the original certificate}\label{finite}

The analytic refinement leaves only \(q<5\cdot10^7\). We retain the
original programs and their recorded search over the larger interval
\(q<10^{10}\), so the verification below covers every remaining case.


\begin{lemma}\label{inert}
If \(h(K)=1\), then every prime \(p<q/4\) is inert in \(K\).
\end{lemma}
\begin{proof}
If \(p\) splits or ramifies, there is an integral ideal of norm \(p\).
It is principal when \(h(K)=1\), so it has a generator
\(\alpha\in\OK\) with \(\norm{\alpha}^2=p\).
A rational integer cannot have prime norm in a quadratic field.
Every nonrational algebraic integer in \(K\) can be written
\((a+b\sqrt{-q})/2\), with integers \(a,b\) and \(b\ne0\);
the allowed parities depend on \(D\).
Consequently,
\[
 p=\norm{\alpha}^2=\frac{a^2+qb^2}{4}\ge\frac q4,
\]
a contradiction.
\end{proof}

For \(q>8\), Lemma~\ref{inert} applies to \(p=2\).
Thus \(2\) is inert, \(D\) is odd, and \(D\equiv5\pmod8\);
equivalently, \(q\equiv3\pmod8\).
The negative fundamental discriminants with \(q\le8\) are
\(-3,-4,-7,-8\).

For every remaining candidate, therefore, it is enough to examine
\[
 11\le q<10^{10},\qquad q\equiv3\pmod8.
\]
For an odd prime \(p<q/4\), exclude \(q\) whenever \(-q\) is zero
or a square modulo \(p\).
The search deliberately includes nonfundamental \(q\), so no
squarefreeness test can accidentally discard a valid candidate.

\subsection*{Exact computation}
Two programs accompany this note:
\begin{itemize}
\item \texttt{verify\_sieve.cpp} builds permitted residue classes using
the Chinese remainder theorem, then checks candidates individually
using tables of squares modulo primes.
\item \texttt{verify\_independent.py} obtains the initial residue
classes by direct scanning and Euler's criterion. It then checks
whole arithmetic progressions using periodic integer bitsets.
\end{itemize}

Both first impose the conditions for \(2,3,5,7,11,13,17\).
The resulting wheel has modulus
\[
 M=8\cdot3\cdot5\cdot7\cdot11\cdot13\cdot17
   =2042040
\]
and \(1440\) permitted residue classes.
For \(437\le q<10^{10}\), every prime at most \(109\) satisfies
\(4p<q\). The programs check the remaining primes from \(19\)
through \(109\), and find no surviving \(q\).
For \(11\le q<437\), they apply only the prime conditions
with \(4p<q\), and find exactly
\[
 11,\quad19,\quad43,\quad67,\quad163.
\]
Both programs retain the original exclusive upper endpoint
\(10^{10}\). The refined analytic argument already excludes every
\(q\ge5\cdot10^7\).
These calculations use exact integer arithmetic throughout.

\subsection*{Checking the surviving fields}
The Python verifier independently enumerates all primitive reduced
positive binary quadratic forms of each remaining discriminant.
The unique reduced form in each case is as follows:
\[
\begin{array}{c|rrrrrrrrr}
q&3&4&7&8&11&19&43&67&163\\ \hline
(a,b,c)&
(1,1,1)&(1,0,1)&(1,1,2)&(1,0,2)&
(1,1,3)&(1,1,5)&(1,1,11)&(1,1,17)&(1,1,41)
\end{array}
\]
In the enumeration, \(\norm b\le a\le c\),
\(b^2-4ac=-q\), \(\gcd(a,b,c)=1\), and \(b\ge0\) on
the boundaries \(\norm b=a\) or \(a=c\).
The standard bound \(a\le\sqrt{q/3}\) makes the enumeration finite.
The classical correspondence with ideal classes verifies
class number one for these nine fields.

\section{Reproduction, scope, and verification}\label{sec-audit}

Run the following commands from this version's directory:
\begin{verbatim}
g++ -O2 -std=c++17 -Wall -Wextra -pedantic \
    verify_sieve.cpp -o verify_sieve
./verify_sieve
python3 verify_independent.py
python3 verify_improved_bound.py
\end{verbatim}
The programs require no number-theory library or external data.
The original sieve outputs are \texttt{sieve\_cpp\_output.txt} and
\texttt{sieve\_python\_output.txt}; they are reused verbatim from version 1.
The two algorithms agree, and reduced-form enumeration checks all nine
survivors. The original Python verifier also checks the rational
inequalities used in the original analytic estimate.

The new certificate, \texttt{verify\_improved\_bound.py}, checks
23 finite inequalities and identities using exact integers and rational
numbers. Its output is \texttt{improved\_bound\_checks.txt}.
For logarithms and exponentials it uses rational Taylor bounds:
if \(x\ge0\), \(n+2>x\), and \(S_n(x)=\sum_{k=0}^n x^k/k!\), then
\[
 S_n(x)\le e^x\le S_n(x)
 +\frac{x^{n+1}}{(n+1)!}\frac1{1-x/(n+2)}.
\]
The tail bound follows by bounding successive term ratios by
\(x/(n+2)<1\). Monotonicity of the exponential converts these to
certificates for the logarithmic inequalities. The standard bounds
\(3<\pi<22/7\) are mathematical inputs.

These numerical checks do not constitute formal verification of the
analytic deductions or of Banks's lemmas. The classical inputs include
Poisson summation for a Gaussian, quadratic ideal theory,
\(\zK=\zeta L\), and reduction of positive binary quadratic forms.
No Lean formalization is supplied. Banks's separate small-conductor
zero computation is not used.

\paragraph{Version-specific review status.}
The repository owner reports that the original version's mathematical
argument was cross-checked with a different frontier model. No detailed
report of that cross-check is bundled here. The refinement in
Sections~\ref{sec-refined-banks}--\ref{sec-refined-theta} has not yet
received a separate-model mathematical review. Its arithmetic certificate
passes, and the typeset paper has been rendered and visually inspected.
The repository's \texttt{REVIEW.md} records the scope of these checks.

\paragraph{AI disclosure and attribution.}
Codex produced the route selection, explicit estimates, proof text,
and all three verification programs. The two sieve algorithms were
written by the same assistant. Their agreement checks computation;
it does not supply an independent review of the analytic proof.
The external determinant results are due to the cited Banks manuscript,
which attributes its underlying argument to the OpenAI manuscript.
This contribution is prepared for Tao's collaborative challenge~\cite{tao}.
It does not claim an independent proof of Banks's results.

The original release and this cumulative version have separate sources,
PDFs, checksums, and fixed version references. Source provenance is
recorded in \texttt{sources/manifest.json}. The inherited provenance
also records a Green manuscript consulted during the original research;
it is not an input to either proof here.

\begin{thebibliography}{9}
\bibitem{banks}
W.~D.~Banks,
\emph{On the OpenAI proof that Landau--Siegel zeros do not exist},
11 October 2026, arXiv:2610.15483v1. Theorem~1.1 for the original
bound; Lemma~4.2, Lemma~5.1, and Corollary~6.3 with their preceding
construction and hypotheses for the refinement.
\url{https://arxiv.org/html/2610.15483v1}.
\bibitem{tao}
T.~Tao,
\emph{An experimental AI-math challenge: using tenth discriminant
bounds to explore arguments ruling out Siegel zeroes},
20 October 2026.
\href{https://terrytao.wordpress.com/2026/10/20/an-experimental-ai-math-challenge-using-tenth-discriminant-bounds-to-explore-arguments-ruling-out-siegel-zeroes/}
{Online challenge}.
\end{thebibliography}
\end{document}
